% 全局排版设置：字体、页边距、公式/图表编号及交叉引用。
\usepackage[a4paper,top=23mm,bottom=23mm,left=23mm,right=23mm]{geometry}
\usepackage{amsmath,amssymb,mathtools,bm}
\usepackage{graphicx,xcolor,tabularx,booktabs,array,calc}
\usepackage{caption,float,placeins,fancyhdr,enumitem,needspace}
\usepackage[unicode,hidelinks]{hyperref}
\usepackage{bookmark,microtype}
\setmainfont{TeX Gyre Termes}
\setsansfont{TeX Gyre Heros}
\setmonofont{Latin Modern Mono}
\input{figures/q1_case}

\linespread{1.18}
\setlength{\parindent}{2em}
\setlength{\parskip}{1.5pt}
\setlength{\emergencystretch}{2em}
\clubpenalty=10000\widowpenalty=10000\displaywidowpenalty=10000
\raggedbottom

% 每问为一个一级节；章节和公式均自动编号，手改时不用维护数字。
\setcounter{secnumdepth}{2}
\ctexset{
  section={name={问题,},number=\chinese{section},format=\Large\bfseries,beforeskip=16pt,afterskip=9pt},
  subsection={format=\large\bfseries,beforeskip=12pt,afterskip=6pt}
}
\numberwithin{equation}{section}
\renewcommand{\theequation}{\thesection-\arabic{equation}}
\newcounter{algorithm}[section]
\renewcommand{\thealgorithm}{\thesection-\arabic{algorithm}}
\setlist[enumerate]{label=\arabic*.,leftmargin=2.2em,itemsep=2pt,topsep=5pt,parsep=0pt}
\setlist[itemize]{leftmargin=2em,itemsep=2pt,topsep=4pt,parsep=0pt}

% 图、表在全篇连续编号。使用 \label 与 \ref，避免写死编号。
\captionsetup{font=small,labelfont=bf,labelsep=quad,skip=6pt,justification=justified,hypcap=false}
\renewcommand{\arraystretch}{1.25}
\setlength{\tabcolsep}{4pt}
\renewcommand{\topfraction}{0.92}
\renewcommand{\bottomfraction}{0.85}
\renewcommand{\textfraction}{0.06}
\renewcommand{\floatpagefraction}{0.65}
\makeatletter
\setlength{\@fptop}{0pt}
\setlength{\@fpsep}{14pt}
\setlength{\@fpbot}{0pt plus 1fil}
\makeatother

\pagestyle{fancy}\fancyhf{}
\fancyhead[C]{\small 无线电干扰源的主动定位与搜索清除}
\fancyfoot[C]{\thepage}
\renewcommand{\headrulewidth}{0.3pt}
\setlength{\headheight}{15pt}
\hypersetup{pdftitle={无线电干扰源的主动定位与搜索清除},pdfauthor={}}
\XeTeXgenerateactualtext=1
